% problem-000085 1 聚合物合成与热力学
\section*{1. 聚合物合成与热力学}
\textit{下列为分问。}

\subsection*{1-1}
画出⼩分⼦ $\mathrm { C _ { 3 } F _ { 4 } }$ 可能的Lewis结构，标出三个碳原⼦成键时所采⽤的杂化轨道。

\subsection*{1-2}
现有铝镁合⾦样品⼀份，称取 2.56 g 样品溶于 HCl，产⽣ 2.80 L H_{2} (101.325 kPa, $0 ~ ^ { \circ } \mathrm { C } )$ ，通过计算，确定此合⾦组成。

\subsection*{1-3}
构筑可循环再⽣的聚合物材料是解决⽬前⽩⾊污染的有效途径之⼀。

\subsection*{1-3}
-1 通常单官能度的单体⽆法参与聚合反应中的链增⻓，但单体A可以与B发⽣反应形成聚合物。画出该聚合物的结构式。

\subsection*{1-3}
-2 单体A、C、D按物质的量⽐例6:3:2进⾏共聚合，得到的聚合物不能进⾏热加⼯和循环利⽤；但若在共聚合时加⼊⼀定量(\~10%)的B，得到的聚合物⼜具备了可热加⼯和循环利⽤的性能。简述原因。

（图：）

A

（图：）

B

（图：）

C

（图：）

D
